\documentclass[12pt]{article}
\usepackage[utf8]{inputenc}
\usepackage[T1]{fontenc}
\usepackage{hyperref}
\title{Dharawal Language Resources for K-12 Curriculum Development}
\author{Compiled by Claude Sonnet 5 (claude-sonnet-5) for Andr\'as Kornai}
\date{}
\begin{document}
\maketitle
\section{Introduction}
Dharawal (ISO 639-3 tbh; Glottocode thur1254; Pama-Nyungan) is a dormant language of the Sydney--Illawarra region of New South Wales that is now being revived. Search results indicate a small number of proficient speakers (reported as 13 in Gujaga Foundation material) and an active, community-led revitalisation effort, so the language is at an early stage of revival. The resources below are only those that appeared in web search results.
\section{Primary Educational Resources}
\begin{itemize}
\item \emph{Dharawal Words, Phrases and Activities}, compiled and edited by Dr Jodi Edwards, illustrated by Merindah Funnell, with input from Elder Aunty Barb Simms; ANSTO funded its printing, along with a Dharawal counting and colouring book. \url{https://www.ansto.gov.au/news/new-language-resource-launched-as-part-of-local-aboriginal-cultural-recognition}
\item Briar Road Public School runs a weekly Dharawal language program delivered with Gujaga Foundation educators. \url{https://briarrd-p.schools.nsw.gov.au/learning-at-our-school/languages.html}
\item Oyster Bay Public School opened an Aboriginal languages room with language lessons for K--2 and Years 3--6. \url{https://www.theleader.com.au/story/8965386/oyster-bay-school-opens-aboriginal-languages-room/}
\end{itemize}
\section{Linguistic and Reference Materials}
\begin{itemize}
\item A UNSW--Gujaga partnership will use archival resources to describe the language with Dharawal knowledge holders and Elders. \url{https://www.unsw.edu.au/newsroom/news/2022/09/reclaiming-language-is-key-to-aboriginal-cultural-identity-and-w}
\item A dictionary compiled by the late Les Bursill was mentioned in coverage of the ANSTO resource; no link or availability was found. \url{https://www.ansto.gov.au/news/new-language-resource-launched-as-part-of-local-aboriginal-cultural-recognition}
\end{itemize}
\section{Religious and Community Resources}
The Gujaga Foundation is described as the peak organisation for language, cultural and research activities in the La Perouse Aboriginal community. The Paul Ramsay Foundation reports a five-year partnership to expand its team of language educators. \url{https://www.paulramsayfoundation.org.au/news-resources/key-aboriginal-language-for-sydney-region-set-to-be-revitalised} \quad \url{https://nit.com.au/05-12-2024/15235/exclusive-dharawal-language-revitalisation-project-given-5m-boost}
\section{Digital and Audio Resources}
A companion website, Dharawal Words, was mentioned in search summaries, but its URL did not appear in the results and is therefore not cited. No audio resources were found.
\section{Curriculum Development Context}
Reports say the Dharawal program runs in more than 50 schools and childcare centres in Sydney (Paul Ramsay Foundation, above). A curriculum could build on the Gujaga Foundation programs, the ANSTO word-and-activity book, and the planned UNSW descriptive work. No full textbook or grammar was found.
\end{document}